\chapter*{Özet}
\setstretch{1.25}
\addcontentsline{toc}{section}{\footnotesize \scshape Özet}

\vspace{-0.6cm}
\rule{\textwidth}{0.3pt}\vspace{-0.35cm}

\begin{justify}
\begin{footnotesize}
Bu çalışmada Strogatz kitabının atanmış 6. bölümünün Türkçeye çevrilmesi ve bölüm sonu sorularının ayrıntılı çözümlerinin hazırlanması amaçlanmaktadır. Belge, enstitü tez formatı esas alınarak düzenlenmiş ve çeviri metni ile çözüm adımlarının sistemli biçimde yerleştirilebilmesi için sade bir iskelet haline getirilmiştir.

\textbf{Anahtar Kelimeler:} \anahtarkelimeler
\end{footnotesize}
\end{justify}

\chapter*{Abstract}
\setstretch{1.25}
\addcontentsline{toc}{section}{\footnotesize \scshape Abstract}

\vspace{-0.6cm}
\rule{\textwidth}{0.3pt}\vspace{-0.35cm}

\begin{justify}
\begin{footnotesize}
This document is prepared as a thesis-format template for translating the assigned Chapter 6 of the Strogatz book into Turkish and for writing detailed solutions to the end-of-chapter problems.

\textbf{Keywords:} \Keywords
\end{footnotesize}
\end{justify}
\newpage
